\chapter{Events, focus, input parsing}
\label{ch:events}

A TUI consumes events: keystrokes, mouse clicks, paste, resize, focus
changes. This chapter dissects how \maya{} parses these from the
ragged byte stream the terminal delivers, and how it routes them to
the right handler.

\section{What the terminal sends}

In raw mode (Chapter~\ref{ch:terminal}) every keystroke arrives as
one or more bytes on \code{stdin}. The mapping is:

\begin{center}
\small
\begin{tabular}{ll}
\textbf{Key} & \textbf{Bytes} \\
\hline
\code{a}                & \code{0x61} \\
\code{Ctrl-A}           & \code{0x01} \\
\code{Esc}              & \code{0x1B} \\
\code{Alt-a}            & \code{0x1B 0x61} \\
\code{Up}               & \code{0x1B [ A} \\
\code{Down}             & \code{0x1B [ B} \\
\code{F1}               & \code{0x1B O P} \\
\code{Shift-Tab}        & \code{0x1B [ Z} \\
\code{Mouse click (SGR)} & \code{0x1B [ < button ; col ; row M} \\
\code{Paste}            & \code{0x1B [ 200~ ... 0x1B [ 201~} \\
\end{tabular}
\end{center}

The parser has to:

\begin{enumerate}[leftmargin=*]
  \item Recognise prefixes incrementally (a single \code{0x1B}
    might be the start of a longer sequence, or it might be a
    lone Escape keypress).
  \item Time out: if no bytes follow \code{0x1B} within a few
    milliseconds, treat it as a lone Esc.
  \item Disambiguate identical prefixes (some terminals send
    \code{ESC[OP} for F1, others \code{ESC[[A}).
  \item Handle bracketed paste --- treat the whole \code{200~...
    201~} block as one event, not a stream of fake keystrokes.
\end{enumerate}

The parser lives in \file{maya/src/terminal/input.cpp}. It is a
state machine driven by incoming bytes; we will not reproduce it
here, but the shape of what it produces is the interesting bit.

\section{The \code{Event} variant}

What the parser \emph{outputs} is a variant:

\begin{lstlisting}
struct Key {
    KeyCode  code;          // Char, Up, Down, F1, ...
    char32_t ch = 0;        // for KeyCode::Char
    Mods     mods{};        // bitset: Ctrl, Alt, Shift
};

struct Mouse {
    int  x = 0, y = 0;
    MouseButton button;     // Left, Middle, Right, ScrollUp, ScrollDown, None
    MouseAction action;     // Press, Release, Move, Drag
    Mods         mods{};
};

struct Resize { int width = 0, height = 0; };
struct Paste  { std::string text; };
struct Focus  { bool gained = true; };

using Event = std::variant<Key, Mouse, Resize, Paste, Focus>;
\end{lstlisting}

Every event \maya{} delivers is one of these. The renderer never sees
raw bytes; the user code never sees raw bytes. The parser is the
choke point.

\subsection{Why types, not strings?}

Some libraries deliver \emph{key chord strings} (\code{"Ctrl+Up"}).
\maya{} delivers structured values. Two reasons:

\begin{itemize}[leftmargin=*]
  \item Pattern matching is exhaustive. \code{std::visit} on
    \code{Event} forces you to handle every alternative.
  \item No string parsing in user code. To check ``user pressed
    Ctrl-Up,'' you read the \code{Key} struct's fields, not a
    string.
\end{itemize}

\section{Event helpers}

For the common pattern of ``did the user press this key,'' \maya{}
provides predicate helpers:

\begin{lstlisting}
bool key(const Event& e, char32_t ch);                   // any key press of ch
bool ctrl(const Event& e, char32_t ch);                  // ctrl-ch
bool alt(const Event& e, char32_t ch);                   // alt-ch
bool key(const Event& e, KeyCode code);                  // arrows, F-keys, ...
bool mouse_clicked(const Event& e, MouseButton btn);     // any click
bool mouse_clicked_in(const Event& e, Rect r, MouseButton b);

template <typename F>
void on(const Event& e, char32_t ch, F&& fn);             // syntactic sugar
\end{lstlisting}

Usage:

\begin{lstlisting}
if (key(ev, 'q'))                  return false;            // quit
if (key(ev, KeyCode::Up))          scroll_up();
if (ctrl(ev, 'c'))                 cancel();
if (mouse_clicked(ev, MouseButton::Left)) select_under_cursor();
\end{lstlisting}

\section{From event to Msg --- \code{key\_map}}

A \code{Program}'s \code{subscribe} function returns a \code{Sub<Msg>}.
The most common subscription is ``map keys to messages,'' and
\maya{} ships a builder for it:

\begin{lstlisting}
template <typename Msg>
Sub<Msg> key_map(std::initializer_list<std::pair<char32_t, Msg>> bindings);
\end{lstlisting}

In use:

\begin{lstlisting}
static Sub<Msg> subscribe(const Model&) {
    return key_map<Msg>({
        {'+',  Inc{}},
        {'-',  Dec{}},
        {'q',  Quit{}},
    });
}
\end{lstlisting}

The runtime, given this \code{Sub}, knows that when the user
presses \code{+} it should construct an \code{Inc\{\}} \code{Msg}
and feed it through \code{update}. Compose more elaborate
subscriptions by combining \code{key\_map}, \code{mouse\_map},
\code{tick} (timers), and \code{batch}.

\section{Focus}

When the user clicks into a different terminal window or tab, the
emulator may send a focus event (\code{ESC [ I} for focus-gained,
\code{ESC [ O} for focus-lost) if focus reporting is enabled. The
parser emits a \code{Focus\{gained: true|false\}} event; \maya{}
programs typically use it to pause animations or dim the UI when
unfocused.

Inside the app, \maya{} has a focus manager
(\file{maya/include/maya/core/focus.hpp}) that tracks which widget
has logical focus, so Tab/Shift-Tab can move focus between input
fields. The widget catalogue in Chapter~\ref{ch:widgets} expects
this manager to exist.

\section{Resize}

When \code{SIGWINCH} fires, the runtime resamples the terminal size
with \code{ioctl(TIOCGWINSZ)} and emits a \code{Resize\{w, h\}}
event. The renderer recomputes layout against the new size and
repaints. Most user code does not need to handle resize explicitly;
the layout engine reflows automatically.

If you do want to react (for example, to re-load a configuration
based on the new aspect ratio), pattern-match on \code{Resize} in
your event handler.

\section{Paste}

Modern terminals support \textbf{bracketed paste}: when the user
pastes a block of text, the terminal wraps the pasted bytes in
\code{ESC[200\textasciitilde} ... \code{ESC[201\textasciitilde}.
\maya{}'s parser collects all bytes between these markers into one
\code{Paste\{text\}} event. This is essential because pasted text
often contains literal newlines and tabs that would otherwise look
like Enter and Tab keypresses --- which would, for example, submit a
chat message after every pasted newline.

To enable bracketed paste, write \code{ESC[?2004h} on startup
(\maya{} does this automatically). To disable, write \code{ESC[?2004l}.

\section{The polling loop}

How does \maya{} actually wait for events? On POSIX, the runtime
uses \code{poll(2)} over three file descriptors:

\begin{itemize}[leftmargin=*]
  \item \code{stdin} --- for key/mouse/paste bytes.
  \item A \code{signalfd}-or-equivalent --- for \code{SIGWINCH}.
  \item A \code{wake\_fd} --- a pipe used to interrupt the poll
    when a background task (a \code{Cmd::Task}) dispatches a
    message.
\end{itemize}

The wake\_fd is a classic Unix trick: when a background thread
wants to feed a \code{Msg} into the event loop, it pushes the
message onto a thread-safe queue and writes one byte to the pipe.
The poll wakes up, drains the pipe, drains the queue, and
processes the message. This is how \maya{}'s pure-functional
\code{update} cooperates with real concurrency.

\section{Windows}

Windows has no \code{poll(2)} for keyboard input. The Windows
implementation uses \code{ReadConsoleInputW} on a dedicated reader
thread, posting events onto a shared queue that the main thread
polls. The interface --- the \code{Event} variant --- is the
same; the plumbing differs.

\begin{recap}
The terminal sends raw bytes; \maya{}'s parser turns them into a
typed \code{Event = variant<Key, Mouse, Resize, Paste, Focus>}.
Helpers (\code{key}, \code{ctrl}, \code{on}) make pattern-matching
ergonomic. The polling loop ties stdin, signals, and a wake-fd
together so that background tasks can interrupt the main loop.
Bracketed paste prevents pasted text from being misread as
keystrokes.
\end{recap}

\section*{Workshop}

\begin{enumerate}[leftmargin=*]
  \item Write a tiny program that prints every \code{Event} it
    receives, then run it and try unusual keys (F-keys, function
    + modifier combos, mouse, paste).
  \item Read the first 200 lines of \file{maya/src/terminal/input.cpp}.
    Identify the state machine's states (waiting for ESC, in CSI,
    in OSC, ...).
  \item Build a small program that uses mouse events to drag a
    coloured square around the screen. The squares' position
    lives in a \code{Signal<Point>}.
\end{enumerate}
